\documentclass[12pt,a4paper]{article}
\usepackage[utf8]{inputenc}
\usepackage[T1]{fontenc}
\usepackage{graphicx}
\usepackage{booktabs}
\usepackage{hyperref}
\usepackage{amsmath}
\usepackage{caption}
\usepackage{float}
\usepackage{geometry}
\usepackage{array}

% \usepackage{tabularx}
\geometry{margin=1in}

\title{Analysis of Battery Bidding Behaviours Through Classification Algorithms in UK Power Markets}

\author{NYDX6\\
Energy Institute, University College London}

\date{}

\begin{document}

\maketitle

\begin{abstract}
This study applies machine learning classification algorithms to analyse bidding behaviours of battery energy storage systems (BESS) across multiple UK electricity markets during 2024. Using comprehensive data from Elexon and NESO covering four market segments---including Enduring Auction Capability (EAC) services, Balancing Mechanism, Balancing Reserve, and wholesale markets---the research investigates whether battery operators possess distinct strategic ``fingerprints'' that differentiate their market participation approaches.

Random Forest and Naive Bayes algorithms were trained to classify battery operators based on their normalized bidding patterns across 48 half-hourly settlement periods. The analysis encompasses bidding data from major battery operators including Tesla Motors, EDF Energy Customers, Statkraft Markets, Limejump Energy, Flexitricity, and others, representing the comprehensive behavioural analysis of UK battery storage market participation. Decision path clustering analysis also reveals tactical sub-patterns within strategic groups.

These findings provide empirical evidence that UK battery storage markets support strategic diversity. The research also establishes a methodological framework for behavioural analysis of BESS in electricity markets, contributing to understanding of how market design influences strategic differentiation among energy storage operators.
\end{abstract}

\noindent\textbf{Keywords:} battery energy storage systems; bidding behaviour; UK power markets; machine learning classification; ancillary services

\section{Summary}
Batteries are playing an increasingly pivotal role in electricity markets worldwide. The global deployment of battery energy storage has surged dramatically in recent years---for example, the United States has seen its installed battery storage capacity surge from around 2.4\,GW in 2020 to exceeding 30\,GW by early 2025, marking over a tenfold increase \cite{eia2024a}. This rapid growth underscores the rising importance of battery energy storage systems (BESS) in enabling a more flexible and low-carbon grid. These systems operate by storing electricity in chemical form and later releasing it when needed, essentially acting as a buffer that absorbs surplus generation and discharges during peak demand. Through this capability, batteries help balance supply and demand in real time and facilitate the integration of intermittent renewable resources into the grid \cite{martins2021}.

The UK serves as an illustrative case of a fast-evolving battery market where companies must adapt quickly. There are multiple revenue opportunities for batteries in the UK---from the Balancing Mechanism (BM) and ancillary frequency response services to reserve markets like Short-Term Operating Reserve (STOR) and the wholesale energy market---and operators often engage in ``revenue stacking'' by participating across several of these markets \cite{bush2024}. This flexibility allows a storage asset to maximize income, but it also means that companies continually adjust their bidding strategies as market rules and competition change. In recent years the UK grid has introduced new ancillary services (for example, a Quick Reserve product launched in late 2024 alongside a Balancing Reserve), which significantly boosted revenues for battery units able to provide rapid response \cite{jennings2025a}. At the same time, increasing competition among storage providers has started to compress profit margins in established services---a phenomenon known as price cannibalization \cite{jennings2025b}. As a result, battery companies in Britain are frequently recalibrating how they bid into each market segment to remain competitive and capture the best value from their assets.

This research utilizes comprehensive bidding data sourced from two key UK electricity market organizations: Elexon and the National Energy System Operator (NESO). Elexon serves as the Settlement Administration Agent for the UK electricity market, and NESO manages the transmission network and procures ancillary services. From publicly available datasets, we collected Balancing Mechanism Unit (BMU) information, detailed bid-offer data from the Balancing Mechanism, Physical Notification, and EAC sell orders. While detailed wholesale electricity market data and capacity market information remain commercially sensitive and are not publicly disclosed, Physical Notification data are employed as a proxy indicator for wholesale market trading positions.

To the best of knowledge, this analysis represents the first comprehensive study of its kind to apply machine learning classification techniques to battery storage bidding behaviours in the UK market, providing novel insights into the strategic bidding patterns of major battery operators using publicly available market data.

Given this dynamic context, three key research questions are investigated:
\begin{enumerate}
    \item Can machine learning classifiers (specifically Random Forest algorithms) successfully identify battery operators from their daily bidding patterns, which indicate that companies possess distinct strategic approaches to market participation?
    \item What distinctive patterns in bidding strategies and market participation behaviours can be identified for different battery operators?
    \item How do different battery operators stack their bidding across multiple markets, and do the identified bidding behaviours correspond to distinct market preferences or varied approaches within the same markets?
\end{enumerate}

This research applies classification algorithms to bidding data from multiple market arenas for a set of large battery operators. By training models to identify which company submitted a given bid, this study tests whether each firm has a unique strategic ``fingerprint'' or whether their bidding patterns demonstrate convergence.

\section{Materials and Methods}

\subsection{Data Collection}
This study analyses battery storage bidding behaviours across multiple UK electricity markets during the period 2024/01/01 to 2024/12/31. All volume-based bidding behaviours are measured in MW/hour, with normalisation method mentioned in next subsection.

Data were collected from publicly available APIs provided by NESO and ELEXON. The data collection strategy focused on capturing bidding decisions rather than revenue outcomes. Table~\ref{tab:datasources} summarises the data sources and key columns used.

\begin{table}[H]
\centering
\caption{Data collection overview}
\label{tab:datasources}
% \begin{tabular}{@{}ll>{\raggedright}p{2.5cm}>{\raggedright}p{3.5cm}@{}}
% \begin{tabular}{@{}p{2.2cm}p{2.5cm}p{5.5cm}p{4.8cm}@{}}
\begin{tabular}{@{}>{\raggedright\arraybackslash}p{2.2cm}
                >{\centering\arraybackslash}p{2.5cm}
                >{\raggedright\arraybackslash}p{5.5cm}
                >{\raggedright\arraybackslash}p{4.8cm}@{}}
\toprule
Market & Data Source & Used Columns & Column Explanations \\ \midrule
EAC & NESO Data Portal & deliveryStart, deliveryEnd, auctionUnit, serviceType, auctionProduct, quantity & Service delivery time windows, BMU identifier, service classification, volume (MW) \\
Balancing Mechanism & ELEXON BMRS Bid-Offer & settlementDate, settlementPeriod, nationalGridBmUnit, levelFrom, levelTo & Trading day and 30-min period, BMU identifier, MW capacity range \\
Wholesale & ELEXON BMRS PN & settlementDate, settlementPeriod, levelFrom, levelTo, nationalGridBmUnit & Planned generation/consumption range (MW) \\
Balancing Reserve & NESO Balancing-Reserve Sell Orders & auctionID, auctionUnit, basketID, serviceType, deliveryStart, deliveryEnd, quantity & Time windows, BMU identifier, MW volume \\
BMU Data & ELEXON BMRS BMU LIST & nationalGridBmUnit, leadPartyName, fuelType, demandCapacity, generationCapacity & Company/operator name, technology classification, technical limits (MW) \\
\bottomrule
\end{tabular}
\end{table}

\subsubsection{Data Cleaning}
All bidding quantities were processed as MW/hour for each settlement period. For EAC data with different temporal structures, volumes were converted to half-hourly equivalents labelled by settlement periods. Battery assets were identified using the \texttt{fuelType} field in Elexon's BMU list, filtering for units marked as \texttt{BATTERY}. However, some battery installations may appear under \texttt{OTHER} fuel type categories; notable examples include Jamesfield 1, Jamesfield 2, Little Raith, and Roaring Hill, which were included after manual verification.

\subsubsection{Data Exclusions}
The Capacity Market was excluded due to lack of publicly available detailed bidding data. Imbalance/cash-out settlements were excluded as they represent settlement outcomes rather than strategic bidding. STOR data were collected but formed a very small dataset with minimal impact. Wholesale market analysis relies on Physical Notifications as a proxy, following the methodology developed by Modo Energy~\cite{modo2024}.

\subsection{Bidding Behaviour Data Generation}
The longitudinal bidding data were transformed into a cross-sectional format where each row represents a single instance (day and BMU) and the 48 settlement periods serve as feature columns, containing normalised bid volumes. Table~\ref{tab:cross_sec} shows an example.

\begin{table}[H]
\centering
\small
\caption{Example of cross-sectional data format for bidding behaviours.}
\label{tab:cross_sec}
\begin{tabular}{@{}llllllll@{}}
\toprule
deliveryDate & auctionUnit & participant & 1 & 2 & \dots & 47 & 48 \\
\midrule
2024-01-01 & AG-GSTK05 & STATKRAFT MARKETS GMBH & 0.0 & 0.0 & \dots & 0.216 & 0.216 \\
2024-01-01 & AG-HLIM04 & LIMEJUMP ENERGY LIMITED & 0.0 & 0.0 & \dots & -0.037 & -0.037 \\
\dots & \dots & \dots & \dots & \dots & \dots & \dots & \dots \\
\bottomrule
\end{tabular}
\end{table}

\subsubsection{EAC Services}
Volume extraction uses the \texttt{quantity} column for each delivery window. The \texttt{auctionProduct} identifier categorises products into high frequency services (DCH, DMH, DRH, NQR: absorb energy) and low frequency services (DCL, DML, DRL, PQR: inject energy). For baskets containing parent and child orders, quantities are summed.

\subsubsection{Balancing Mechanism}
Data preprocessing excludes rows where both \texttt{levelFrom} and \texttt{levelTo} equal zero. Volume is calculated as the average MW: $(\text{levelFrom} + \text{levelTo}) / 2$. Following BM conventions, bids are associated with reducing output, offers with increasing output.

\subsubsection{Balancing Reserve}
Similar to EAC, quantities are summed across parent-child basket structures. Positive Balancing Reserve (PBR) quantities are divided by generation capacity and Negative Balancing Reserve (NBR) by demand capacity for normalisation.

\subsubsection{Wholesale Market (Physical Notification)}
Processing mirrors the Balancing Mechanism methodology, calculating average MW as $(\text{levelFrom} + \text{levelTo}) / 2$ for each settlement period.

\subsection{Data Normalisation}
Raw volumes are normalised using BMU capacity data to focus on bidding strategy rather than asset scale:
\[
\text{Normalised Volume} = \frac{\text{Raw Volume (MW/h)}}{\text{BMU Capacity (MW/h)}}.
\]
Capacity reference data come from Elexon's BMU dataset, representing generation capacity (GC) and demand capacity (DC). The distributions of these capacities are shown in Figures~\ref{fig:gen_cap_dist} and~\ref{fig:dem_cap_dist}.

\begin{figure}[H]
    \centering
    \includegraphics[width=0.6\textwidth]{media/image10.png}
    \caption{Distribution of generation capacity of BMU.}
    \label{fig:gen_cap_dist}
\end{figure}

\begin{figure}[H]
    \centering
    \includegraphics[width=0.6\textwidth]{media/image11.png}
    \caption{Distribution of demand capacity of BMU.}
    \label{fig:dem_cap_dist}
\end{figure}
\subsection{Data Exploratory Analysis}
Before building classification models, we performed an exploratory analysis to understand the overall participation patterns, temporal distributions, and correlations across market segments.

\subsubsection{EAC Market Participation}
Figure~\ref{fig:eac_monthly} shows the order counts per month across EAC auction products throughout 2024. A dramatic spike in Positive Quick Reserve (PQR) and Negative Quick Reserve (NQR) is observed in December 2024, indicating immediate adoption of the newly launched Quick Reserve service. The correlation heatmap of EAC bidding volumes (Figure~\ref{fig:eac_corr}) reveals strong positive correlations between complementary direction pairs, such as DCH--DCL and NQR--PQR, suggesting that operators bid similar volumes for positive and negative directions within the same service. Strong negative correlations appear between Dynamic Regulation and Dynamic Containment, reflecting different operational roles.

\begin{figure}[H]
    \centering
    \includegraphics[width=0.85\textwidth]{media/image12.png}
    \caption{Monthly bidding order numbers by auction product in EAC.}
    \label{fig:eac_monthly}
\end{figure}

\begin{figure}[H]
    \centering
    \includegraphics[width=0.5\textwidth]{media/image13.png}
    \caption{Correlation matrix of bidding volumes by EAC auction product.}
    \label{fig:eac_corr}
\end{figure}

The distributions of normalised bidding volumes for EAC services are shown together in Figure~\ref{fig:eac_dist}. Low- and high-frequency products are plotted in paired panels, so DCL and DCH, DML and DMH, DRL and DRH, and PQR and NQR can be compared directly within the same graph. Quick Reserve services (NQR, PQR) and Dynamic Moderation (DMH, DML) display highly concentrated distributions with most bids clustered at small volumes. Dynamic Regulation shows stronger asymmetry between high and low products, while Dynamic Containment (DCH, DCL) distributions are broader, indicating more flexible capacity allocation.

\begin{figure}[H]
    \centering
    \includegraphics[width=\textwidth]{media/image14.png}
    \caption{Normalised EAC bidding volume distributions by paired low- and high-frequency products.}
    \label{fig:eac_dist}
\end{figure}

Among the major participants, Tesla leads dynamic‑service bidding, EDF is the strongest broad follower, Statkraft maintains significant presence in DC and QR, and Quick Reserve volumes become visible only at year‑end.

\subsubsection{Balancing Mechanism Participation}
Balancing Mechanism participant dominance is stable across the year. Statkraft remains the largest contributor in most months, followed by EDF, Tesla, and Limejump. Bid and offer count breakdowns for selected operators (Figures~\ref{fig:bm_edf} and~\ref{fig:bm_statkraft}) show that some companies (e.g., Statkraft) are generation‑oriented, while others (e.g., EDF) exhibit more balanced or demand‑side participation.

\begin{figure}[H]
    \centering
    \includegraphics[width=0.7\textwidth]{media/image22.png}
    \caption{BM Bid and offer counts per month for EDF Customers Limited batteries.}
    \label{fig:bm_edf}
\end{figure}

\begin{figure}[H]
    \centering
    \includegraphics[width=0.68\textwidth]{media/image23.png}
    \caption{BM Bid and offer counts per month for Statkraft Markets GmbH batteries.}
    \label{fig:bm_statkraft}
\end{figure}

The correlation matrix of BM decisions across settlement periods (Figure~\ref{fig:bm_corr}) reveals that bidding behaviours cluster by operational time blocks rather than only by adjacent periods: morning periods are highly correlated with other morning periods, and evening periods are correlated with each other. A noticeable correlation also appears between noon and midnight/overnight periods.

\begin{figure}[H]
    \centering
    \includegraphics[width=0.55\textwidth]{media/image21.png} % Note: The original doc had a BM correlation matrix, possibly image25; we use the correct file name. Assume image25.png.
    \caption{Balancing Mechanism settlement period correlation matrix.}
    \label{fig:bm_corr}
\end{figure}

% The normalised volume distribution is strongly concentrated around zero (Figure~\ref{fig:bm_vol_dist}). Temporally, BM bidding activity grew steadily through 2024, reaching a record in December (Figure~\ref{fig:bm_monthly}), likely reflecting shifting revenue opportunities as frequency response margins compressed.

% \begin{figure}[H]
%     \centering
%     \includegraphics[width=0.7\textwidth]{media/image26.png}
%     \caption{Distribution of bid/offer volumes in the Balancing Mechanism.}
%     \label{fig:bm_vol_dist}
% \end{figure}

% \begin{figure}[H]
%     \centering
%     \includegraphics[width=0.8\textwidth]{media/image27.png}
%     \caption{Monthly total bidding count in the Balancing Mechanism.}
%     \label{fig:bm_monthly}
% \end{figure}

\subsubsection{Balancing Reserve Participation}
Balancing Reserve (BR) volume distribution shows a strong central peak at zero with secondary peaks at 0.2 and 0.4 normalised volume (Figure~\ref{fig:eac_dist}), suggesting that certain operators consistently bid at specific fractional capacities. Positive Balancing Reserve (PBR) dominating (Figure~\ref{fig:eac_dist}). Tesla appears the major operator (Figure~\ref{fig:br_operator}).
% \begin{figure}[H]
%     \centering
%     \includegraphics[width=0.7\textwidth]{media/image28.png}
%     \caption{Distribution of bidding volumes in Balancing Reserve.}
%     \label{fig:br_vol_dist}
% \end{figure}

% \begin{figure}[H]
%     \centering
%     \includegraphics[width=0.5\textwidth]{media/image29.png}
%     \caption{Monthly Balancing Reserve bids: PBR vs NBR.}
%     \label{fig:br_monthly}
% \end{figure}

\begin{figure}[H]
    \centering
    \includegraphics[width=0.6\textwidth]{media/image28.png}
    \caption{BR operators' participation by month.}
    \label{fig:br_operator}
\end{figure}

\subsubsection{Wholesale Market (Physical Notification)}
Physical Notifications (PN), used as a proxy for wholesale trading, exhibit a near‑normal distribution centred around zero with a slight positive skew (Figure~\ref{fig:eac_dist}), indicating that batteries more frequently position for export (discharge) than import. Table~\ref{tab:pn_counts} shows strong PN participation across all major operators.

\begin{table}[H]
\centering
\caption{Number of bidding instances in Physical Notification by company (2024).}
\label{tab:pn_counts}
\begin{tabular}{@{}lr@{}}
\toprule
Company & Number of PN submissions \\
\midrule
Statkraft Markets GmbH & 6173 \\
Tesla Motors Limited & 5723 \\
EDF Energy Customers Limited & 5400 \\
Flexitricity Limited & 3106 \\
Limejump Energy Limited & 2484 \\
\bottomrule
\end{tabular}
\end{table}



\subsubsection{Overall Market Participation}
Combining all markets, we observe that Dynamic Containment and Dynamic Moderation tend to be more positive during early/midday blocks and more negative around the evening peak (especially SP31--44). Dynamic Regulation is mostly negative across the day, with occasional positive blocks overnight or late day, suggesting it is often used for charging or state‑of‑charge management. Physical Notification is frequently negative in the morning, strongly positive around SP35--40, and moderately positive near midday, consistent with evening peak discharge and intraday repositioning. Quick Reserve only appears in Month 12 and has a strong structured profile: negative at the start/end of the day, positive across most daytime periods. Balancing Mechanism is mostly mildly negative across all months and periods, appearing as a persistent charging or counter‑positioning layer.

Figure~\ref{fig:overall_heatmap} (a composite visualisation) summarises the average normalised volume per market and settlement period across all operators. This multi‑market view confirms that different services fulfil distinct roles in the operators’ revenue‑stacking strategies.

\begin{figure}[H]
    \centering
    % This figure was not numbered distinctly in the original, we assume image43 or similar.
    \includegraphics[width=\textwidth]{media/image43.png}
    \caption{Overall market participation profiles: average normalised volume per settlement period across all operators.}
    \label{fig:overall_heatmap}
\end{figure}

\subsection{Classification Approach}
Two machine learning algorithms were compared: Naive Bayes as a baseline and Random Forest as the primary classifier. Naive Bayes assumes feature independence, providing a benchmark. Random Forest captures non-linear feature interactions and handles high-dimensional data without requiring explicit independence assumptions.

\subsection{Evaluation Approach}
Given severe class imbalance in the UK battery market, model performance is evaluated using Overall Accuracy, Weighted F1-score, and Macro F1-score. The Macro F1-score is the primary metric as it treats all market participants equally, validating the algorithm's ability to capture diverse strategic fingerprints of smaller operators.

\section{Results}

\subsection{Classification Report}
Table~\ref{tab:class_report} summarises the performance of both classifiers across the seven market segments. Random Forest substantially outperforms Naive Bayes in all markets, indicating that bidding behaviours involve complex feature interactions.

% \begin{table}[H]
% \centering
% \caption{Classification performance summary}
% \label{tab:class_report}
% \begin{tabular}{@{}lcccc@{}}
% \toprule
%  & \multicolumn{2}{c}{Random Forest} & \multicolumn{2}{c}{Naive Bayes} \\
% \cmidrule(lr){2-3} \cmidrule(lr){4-5}
% Market & Accuracy & Macro F1 & Accuracy & Macro F1 \\
% \midrule
% Dynamic Containment & 0.87 & 0.78 & 0.29 & 0.24 \\
% Dynamic Moderation & 0.91 & 0.79 & 0.29 & 0.22 \\
% Dynamic Regulation & 0.95 & 0.93 & 0.41 & 0.32 \\
% Quick Reserve & 0.94 & 0.87 & 0.85 & 0.72 \\
% Balancing Reserve & 0.97 & 0.69 & 0.67 & 0.38 \\
% Balancing Mechanism & 0.77 & 0.64 & 0.30 & 0.24 \\
% Physical Notification & 0.67 & 0.48 & 0.20 & 0.15 \\
% \bottomrule
% \end{tabular}
% \end{table}

\begin{table}[H]
\centering
\caption{Classification performance summary (Accuracy, Macro F1, Weighted F1)}
\label{tab:class_report}
\begin{tabular}{@{}lcccccc@{}}
\toprule
& \multicolumn{3}{c}{Random Forest} & \multicolumn{3}{c}{Naive Bayes} \\
\cmidrule(lr){2-4} \cmidrule(lr){5-7}
Market & Acc & MF1 & WF1 & Acc & MF1 & WF1 \\
\midrule
Dynamic Containment     & 0.87 & 0.78 & 0.87 & 0.29 & 0.24 & 0.29 \\
Dynamic Moderation      & 0.91 & 0.79 & 0.90 & 0.29 & 0.22 & 0.30 \\
Dynamic Regulation      & 0.95 & 0.93 & 0.95 & 0.41 & 0.32 & 0.44 \\
Quick Reserve           & 0.94 & 0.87 & 0.94 & 0.85 & 0.72 & 0.86 \\
Balancing Reserve       & 0.97 & 0.69 & 0.97 & 0.67 & 0.38 & 0.74 \\
Balancing Mechanism     & 0.77 & 0.64 & 0.76 & 0.30 & 0.24 & 0.32 \\
Physical Notification   & 0.67 & 0.48 & 0.64 & 0.20 & 0.15 & 0.17 \\
\bottomrule
\end{tabular}
\end{table}

\subsection{Naive Bayes Baseline}
Confusion matrix examples. The Naive Bayes confusion matrices (Figures~\ref{fig:nb_dc} and~\ref{fig:nb_dm}) reveal poor performance, with most predictions dominated by the majority class.

\begin{figure}[H]
    \centering
    \includegraphics[width=0.7\textwidth]{media/image32.png}
    \caption{Naive Bayes classification confusion matrix -- Dynamic Containment.}
    \label{fig:nb_dc}
\end{figure}

\begin{figure}[H]
    \centering
    \includegraphics[width=0.7\textwidth]{media/image33.png}
    \caption{Naive Bayes classification confusion matrix -- Dynamic Moderation.}
    \label{fig:nb_dm}
\end{figure}

\subsection{Random Forest Performance}
Confusion matrix examples. Random Forest confusion matrices (Figures~\ref{fig:rf_dc} and~\ref{fig:rf_dm}) show much stronger diagonal dominance, demonstrating successful identification of operator-specific strategies.

\begin{figure}[H]
    \centering
    \includegraphics[width=0.7\textwidth]{media/image39.png}
    \caption{Random Forest classification confusion matrix -- Dynamic Containment.}
    \label{fig:rf_dc}
\end{figure}

\begin{figure}[H]
    \centering
    \includegraphics[width=0.7\textwidth]{media/image40.png}
    \caption{Random Forest classification confusion matrix -- Dynamic Moderation.}
    \label{fig:rf_dm}
\end{figure}

\subsection{Random Forest Feature Importance}
Feature importance analysis (Figures~\ref{fig:fi_dc}--\ref{fig:fi_br}) reveals distinct temporal patterns that differentiate operator strategies. Quick Reserve and Dynamic Regulation show extreme concentration in period 48, while Balancing Mechanism discrimination occurs primarily in early periods (1--3). Several services concentrate discriminative power during traditional peak demand periods (17--22, 35--45), suggesting that strategic differentiation is most pronounced when market values are highest.

\begin{figure}[H]
    \centering
    % \captionsetup{font=footnotesize}
    \captionsetup{font=footnotesize,justification=raggedright}
    \begin{minipage}{0.32\textwidth}
        \centering
        \small
        \includegraphics[width=\textwidth]{media/image46.png}
        \caption{DC.}
        \label{fig:fi_dc}
    \end{minipage}
    \begin{minipage}{0.32\textwidth}
        \centering
        \small
        \includegraphics[width=\textwidth]{media/image47.png}
        \caption{DM.}
        \label{fig:fi_dm}
    \end{minipage}
    \begin{minipage}{0.32\textwidth}
        \centering
        \small
        \includegraphics[width=\textwidth]{media/image49.png}
        \caption{QR.}
        \label{fig:fi_qr}
    \end{minipage}
    % \vspace{0.2cm}
    \begin{minipage}{0.32\textwidth}
        \centering
        \small
        \includegraphics[width=\textwidth]{media/image50.png}
        \caption{PN.}
        \label{fig:fi_pn}
    \end{minipage}
    % \hfill
    \begin{minipage}{0.32\textwidth}
        \centering
        \small
        \includegraphics[width=\textwidth]{media/image51.png}
        \caption{BM.}
        \label{fig:fi_bm}
    \end{minipage}
    % \hfill
    \begin{minipage}{0.32\textwidth}
        \centering
        \small
        \includegraphics[width=\textwidth]{media/image52.png}
        \caption{BR.}
        \label{fig:fi_br}
    % \vspace{0.2cm}
    \end{minipage}
\end{figure}

\vspace{1cm}
\subsection{Pattern Clustering}
Decision path clustering was applied within each operator group to identify tactical sub-patterns. The method are chosed from k-mean, DBSCAN and spectral clustering. Clusters were evaluated using silhouette score, and the best clustering solution was selected. The Cramér\’s V association between cluster membership and auction unit/month was computed to estimate whether strategies are asset-specific or temporally driven. But the limitation of small sample size and mixed effect of variables should be awared.


\subsection{Behaviour Patterns by Markets}
Detailed bidding patterns for each market are presented in Figures~\ref{fig:dc_heatmaps}--\ref{fig:br_heatmaps} and discussed below.

\subsubsection{Dynamic Containment}
DC patterns align with Electricity Forward Agreement (EFA) blocks. Two primary charging windows appear: morning (07:00--11:00, EFA block 3) and afternoon/evening (15:00--22:00, blocks 5 and 6). Most companies simultaneously bid for high and low frequency services using loop basket structures. Flexitricity might demonstrates asset-specific strategies (Cram\'er's V = 0.87 with auction unit).

\begin{figure}[H]
    \centering
    \includegraphics[width=0.9\textwidth]{media/image54.png}
    % \includegraphics[width=\textwidth]{media/image55.png}
    % \includegraphics[width=\textwidth]{media/image56.png}
    \caption{Heatmaps of Dynamic Containment bidding patterns by companies.}
    \label{fig:dc_heatmaps}
\end{figure}

\begin{figure}[H]
    \centering
    \includegraphics[width=0.8\textwidth]{media/image68.png}
    \caption{Flexitricity Limited bidding patterns in Dynamic Containment (example of asset-specific strategy).}
    \label{fig:flex_dc}
\end{figure}

\vspace{2cm}
\subsubsection{Dynamic Moderation}
DM shows similar peak-period positioning. EDF, Flexitricity and Statkraft maintain unidirectional strategies; Limejump exhibits contrasting positive/negative bidding, while Tesla balances participation moderately. Cram\'er's V values might indicate temporal factors.(Cram\'er's V = 0.24 with auction unit, 0.41 with month)

\begin{figure}[H]
    \centering
    \includegraphics[width=0.9\textwidth]{media/image80.png}
    \caption{Heatmaps of Dynamic Moderation bidding patterns.}
    \label{fig:dr_heatmaps}
\end{figure}

\vspace{4cm}
\subsubsection{Dynamic Regulation}
DR patterns are remarkably uniform across companies because of the stringent 60-minute minimum energy requirement favouring conservative volume approaches. 

\begin{figure}[H]
    \centering
    \includegraphics[width=0.9\textwidth]{media/image82.png}
    \caption{Heatmaps of Dynamic Regulation bidding patterns.}
    \label{fig:dr_heatmaps}
\end{figure}

% \begin{figure}[H]
%     \centering
%     \includegraphics[width=0.8\textwidth]{media/image89.png}
%     \caption{Tesla Motors Limited bidding patterns in Dynamic Regulation.}
%     \label{fig:tesla_dr}
% \end{figure}

The contrasting temporal behaviours observed in Dynamic Containment and Dynamic Regulation can be directly linked to their distinct technical requirements (Table~\ref{tab:eac_tech_specs}). 
Dynamic Containment’s short 15‑minute duration and post‑fault activation allow operators to aggressively reposition their assets between charging and discharging across settlement periods, resulting in the pronounced intra‑day shifts (e.g.\ noon charging, morning discharge) that align with EFA block price dynamics.
In contrast, Dynamic Regulation mandates continuous provision for 60 minutes at full power over $\pm 0.2$\,Hz; this forces operators to maintain a stable state of charge and a conservative, almost unidirectional bidding profile to avoid energy depletion or saturation within the tendered hour. 
\vspace{0.7cm}
\begin{table}[H]
\centering
\caption{EAC service technical requirements overview}
\label{tab:eac_tech_specs}
\begin{tabular}{@{}>{\raggedright\arraybackslash}p{4.3cm}>{\raggedright\arraybackslash}p{2.3cm}>{\raggedright\arraybackslash}p{2.3cm}>{\raggedright\arraybackslash}p{2.3cm}>{\raggedright\arraybackslash}p{3cm}@{}}
\toprule
Service & Duration Capability & Response Time & Frequency Range & Service Type \\
\midrule
Dynamic Containment & 15 minutes & Sub‑second & $>\pm 0.2$\,Hz & Post‑fault \\
Dynamic Moderation & 30 minutes & Fast‑acting & $\pm 0.1$--$0.2$\,Hz & Pre‑fault (volatile) \\
Dynamic Regulation & 60 minutes & Proportional & Full power at $\pm 0.2$\,Hz & Pre‑fault (continuous) \\
Quick Reserve & 5 minutes & 1 min to full delivery & N/A & Reserve service \\
\bottomrule
\end{tabular}
\end{table}


\vspace{4cm}
\subsubsection{Quick Reserve}
QR strategies diversify significantly. Statkraft bids for charging at night and discharge during daytime; J Aron \& Company shows pronounced discharge; Limejump focuses on reserve provision, and Arenko CleanTech exhibits large volume bids.

\begin{figure}[H]
    \centering
    \includegraphics[width=0.9\textwidth]{media/image91.png}
    \caption{Quick Reserve heatmaps by company.}
    \label{fig:qr_heatmaps}
\end{figure}

\begin{figure}[H]
    \centering
    \includegraphics[width=0.8\textwidth]{media/image98.png}
    \caption{Statkraft Markets GMBH bidding patterns in Quick Reserve.}
    \label{fig:statkraft_qr}
\end{figure}

\vspace{3cm}
\subsubsection{Physical Notification (Wholesale)}
PN patterns indicate peak-hour discharge consistent with energy arbitrage, with more bounded and volatile profiles compared to ancillary services. EDF shows discrete, volatile bidding across periods.

\begin{figure}[H]
    \centering
    \includegraphics[width=0.9\textwidth]{media/image112.png}
    \caption{Physical Notification heatmaps.}
    \label{fig:br_heatmaps}
\end{figure}

\begin{figure}[H]
    \centering
    \includegraphics[width=0.8\textwidth]{media/image111.png}
    \caption{EDF Energy Customers Limited bidding patterns in Physical Notification.}
    \label{fig:edf_pn}
\end{figure}

\vspace{2cm}
\subsubsection{Balancing Mechanism}
In the BM, companies like Centrica and Yuso demonstrate peak-period charging or generation reduction, contrasting with wholesale patterns. Statkraft maintains conventional generation increases; Limejump exhibits heavy unidirectional charging.

\begin{figure}[H]
    \centering
    \includegraphics[width=0.9\textwidth]{media/image121.png}
    \caption{Balancing Mechanism heatmaps.}
    \label{fig:limejump_bm}
\end{figure}

% \begin{figure}[H]
%     \centering
%     \includegraphics[width=0.8\textwidth]{media/image120.png}
%     \caption{Limejump Energy Limited bidding patterns in Balancing Mechanism.}
%     \label{fig:limejump_bm}
% \end{figure}

\begin{figure}[H]
    \centering
    \includegraphics[width=0.8\textwidth]{media/image125.png}
    \caption{EDF Energy Customers Limited bidding patterns in Balancing Mechanism.}
    \label{fig:edf_bm}
\end{figure}


\subsubsection{Balancing Reserve}
BR patterns feature gradual volume ramps and marked temporal differentiation. EDF's discrete bidding behaviour contrasts with other participants, possibly reflecting portfolio management to minimise penalties.

\begin{figure}[H]
    \centering
    \includegraphics[width=0.9\textwidth]{media/image123.png}
    % \includegraphics[width=\textwidth]{media/image124.png}
    \caption{Balancing Reserve heatmaps.}
    \label{fig:br_heatmaps}
\end{figure}
\begin{figure}[H]
    \centering
    % \includegraphics[width=\textwidth]{media/image123.png}
    \includegraphics[width=0.85\textwidth]{media/image124.png}
    \caption{Tesla Motors Limited bidding patterns in Balancing Reserve.}
    \label{fig:br_heatmaps}
\end{figure}

\subsection{Behaviour Pattern Comparison by Companies}

Rather than analysing bidding behaviours in isolated markets, the algorithm tracked the cluster assignments of all battery units simultaneously across market segments. 

\begin{figure}[H]
    \centering
    \includegraphics[width=0.82\textwidth]{media/tesla_consistent_table.png}
    \caption{Tesla Motors Limited cross‑market bidding patterns consistent groups.}
    \label{fig:tesla_cross}
\end{figure}

\textbf{Tesla} uses Dynamic Containment/Moderation/Regulation actively for charging/SOC management and relies on Physical Notification for evening peak discharge; Balancing Reserve plays a supporting role.

\textbf{Statkraft} shows a QR/BR-led discharge strategy, with DR acting as a charging layer; PN is opportunistic around the evening peak.

\textbf{EDF} is more BR/PN-led, with DR consistently negative and DC/DM providing complementary charging.

\textbf{Limejump} displays high-volume dual-directional bidding, while \textbf{Flexitricity} demonstrates asset-specific specialisation across frequency response services, confirmed by consistently high Cram\'er's V values for auction units in all Dynamic services.

\begin{figure}[H]
    \centering
    \includegraphics[width=\textwidth]{media/tesla_subplots.png}
    \caption{Tesla Motors Limited: subplots of individual battery asset bidding patterns across markets. All Tesla assets display highly consistent bidding behaviours, reflecting the company's uniform fleet‑wide strategy.}
    \label{fig:tesla_subplots}
\end{figure}

\begin{figure}[H]
    \centering
    \includegraphics[width=\textwidth]{media/tesla_crossmarket.png}
    \caption{Tesla Motors Limited: cross‑market bidding patterns. Tesla uses DC/DM/DR for charging and SOC management, PN for evening peak discharge, and BR as a supporting layer.}
    \label{fig:tesla_cross}
\end{figure}

\begin{figure}[H]
    \centering
    \includegraphics[width=\textwidth]{media/statkraft_crossmarket.png}
    \caption{Statkraft Markets GmbH: cross‑market bidding patterns. Discharge is led by Quick Reserve and Balancing Reserve, with DR serving as a charging layer and PN providing opportunistic evening peak discharge.}
    \label{fig:statkraft_cross}
\end{figure}


\paragraph{EDF Energy Customers consistent group.}
Two such groups within EDF Energy Customers Limited are particularly illustrative.

\begin{figure}[H]
    \centering
    \includegraphics[width=0.75\textwidth]{media/edf_consistent_table.png}
    \caption{EDF Energy Customers Limited cross‑market bidding patterns consistent groups.}
    \label{fig:tesla_cross}
\end{figure}

\paragraph{Group 1 EDF: BR‑led strategy.}
This cohort (Table~\ref{tab:edf_group1}) is characterised by a strong Balancing Reserve presence. Asset COVNB‑1, for instance, exhibits clearly positive BR blocks around SP6--12 and SP29--36, accompanied by a Physical Notification rise at SP35--40, while Dynamic Regulation remains consistently negative, serving as a charging.

\begin{figure}[H]
    \centering
    \includegraphics[width=0.95\textwidth]{media/edf_crossmarket_0.png}
    \caption{EDF Energy Customers Limited: cross‑market bidding patterns. EDF relies on BR and PN for discharge, with QR in night and DR offering complementary charging.}
    \label{fig:edf_cross}
\end{figure}

\begin{table}[H]
\centering
\caption{EDF Group 1 assets (BR‑led strategy)}
\label{tab:edf_group1}
\begin{tabular}{@{}lcc@{}}
\toprule
\textbf{nationalGridBmUnit} & \textbf{generationCapacity (MW)} & \textbf{gspGroupName} \\
\midrule
AG‑PEDF01 & 49.70 & North Scotland \\
NURSB‑1   & 41.69 & --- \\
\bottomrule
\end{tabular}
\end{table}

\paragraph{Group 3 EDF: DC/DM‑driven dynamic‑service strategy.}
In contrast, Group 3 (Table~\ref{tab:edf_group3}) assets such display positive DC/DM overnight and late day, turning negative during mid‑day transition periods, and DR consistant charging. This group relies on frequency response services rather than balancing services for its revenue.

\begin{figure}[H]
    \centering
    \includegraphics[width=0.95\textwidth]{media/edf_crossmarket.png}
    \caption{EDF Energy Customers Limited: cross‑market bidding patterns. DR consistently negative.}
    \label{fig:edf_cross}
\end{figure}
\vspace{2cm}
\begin{table}[H]
\centering
\caption{EDF Group 3 assets (DC/DM‑driven strategy)}
\label{tab:edf_group3}
\begin{tabular}{@{}lc@{}}
\toprule
\textbf{nationalGridBmUnit} & \textbf{generationCapacity (MW)} \\
\midrule
PINFB‑2 & 26.0 \\
PINFB‑3 & 26.0 \\
PINFB‑4 & 26.0 \\
\bottomrule
\end{tabular}
\end{table}

\vspace{1.5cm}
The EDF boxplot shows that group capacity profiles differ: This suggests that asset size may partly influence group membership, although the evidence is limited because several groups contain only a small number of observations.


\begin{figure}[H]
    \centering
    \includegraphics[width=\textwidth]{media/edf_boxplot.png}
    \caption{EDF generation-capacity distribution by consistent bidding group.}
    \label{fig:edf_box}
\end{figure}

\vspace{8cm}
\textbf{Limejump} displays high-volume dual-directional bidding, while \textbf{Flexitricity} demonstrates asset-specific specialisation across frequency response services, confirmed by consistently high Cram\'er's V values for auction units in all Dynamic services.

\begin{figure}[H]
    \centering
    \includegraphics[width=\textwidth]{media/limejump_crossmarket.png}
    \caption{Limejump Energy Limited: cross‑market bidding patterns.}
    \label{fig:limejump_cross}
\end{figure}

\begin{figure}[H]
    \centering
    \includegraphics[width=\textwidth]{media/flexitricity_crossmarket.png}
    \caption{Flexitricity Limited: cross‑market bidding patterns.}
    \label{fig:flex_cross}
\end{figure}

  Across Dynamic Containment, Dynamic Moderation and Dynamic Regulation, the association with Auction Unit remains consistently above 0.87, while the association with Month is notably lower. This implies that Flexitricity’s bidding strategies are predominantly asset‑driven rather than seasonally adjusted; each battery unit follows a customised strategy likely tailored to its operational characteristics or locational advantages, rather than applying a uniform portfolio‑wide approach.

\begin{table}[H]
\centering
\caption{Cramér’s V between EAC market behavioural clusters and Auction Unit / Month for Flexitricity Limited}
\label{tab:flex_cramers}
\begin{tabular}{@{}lccc@{}}
\toprule
& \multicolumn{3}{c}{Cramér’s V} \\
\cmidrule(lr){2-4}
& Dynamic Containment & Dynamic Moderation & Dynamic Regulation \\
\midrule
Auction Unit & 0.8688 & 0.9115 & 0.8710 \\
Month        & 0.2668 & 0.5293 & 0.2782 \\
\bottomrule
\end{tabular}
\end{table}

As described in the main text, these profiles confirm the existence of distinct strategic fingerprints. Tesla’s diversified portfolio, Statkraft’s QR/BR‑led discharge, EDF’s BR/PN focus, Limejump’s high‑volume dual‑directionality, and Flexitricity’s asset‑specific flexibility together demonstrate that multiple viable business models coexist in the UK battery storage market.


\section{Discussion}
The results demonstrate that machine learning classification can successfully identify distinct strategic fingerprints among UK battery storage operators. Random Forest achieves high macro F1-scores across markets, confirming that operator behaviours involve complex, non-linear interactions that Naive Bayes cannot capture.

\textbf{Research Question 1}: The high classification accuracy of Random Forest, especially in Dynamic Regulation (Macro F1 = 0.93) and Quick Reserve (0.87), provides strong evidence that battery operators possess unique, distinguishable bidding patterns. The drastic performance gap between Random Forest and Naive Bayes highlights that strategic differentiation relies on coordinated temporal patterns rather than isolated settlement-period decisions.

\textbf{Research Question 2}: Several distinctive patterns emerge. Dynamic Containment and Moderation strategies align closely with EFA block structures, with peak charging during morning and evening price peaks. Dynamic Regulation forces strategy convergence due to the 60-minute energy requirement, limiting differentiation. Quick Reserve, as a new market, already shows diverse adoption: Statkraft's day/night split, J Aron's discharge dominance, and Arenko's large-volume bids. In the Balancing Mechanism, some operators (Centrica, Yuso) charge during peak periods, opposing wholesale arbitrage behaviour, reflecting the real-time balancing function.

\textbf{Research Question 3}: Cross-market stacking reveals clear strategic archetypes. Tesla adopts a balanced portfolio, distributing charging/discharging across multiple services with a distinct evening peak discharge via PN. Statkraft and EDF concentrate discharge in BM, BR and QR, using DR primarily for charging. Limejump employs high-volume contrasting bids, suggesting an opportunity-focused approach. Flexitricity's highly asset-specific strategies indicate a flexibility-management business model that tailors each asset to local market conditions. These archetypes suggest that multiple viable business models coexist in the UK battery market.

\vspace{2cm}
\section{Conclusions}
This study successfully applies machine learning classification to UK battery storage bidding behaviours, demonstrating that operators possess distinct strategic fingerprints. Key findings include:
\begin{itemize}
    \item Random Forest classifiers achieve high performance (Macro F1 up to 0.93) in identifying operators, confirming strategic diversity.
    \item Strategic patterns align with market structures: EFA blocks for DC/DM, regulatory constraints for DR, and distinct discharge/charge roles across BM, BR and wholesale.
    \item Four strategic archetypes emerge: portfolio diversification (Tesla), concentrated peak discharge (Statkraft/EDF), high-volume contrasting bids (Limejump), and asset-specific flexibility (Flexitricity).
    \item Feature importance reveals that competitive differentiation peaks during high-value settlement periods.
\end{itemize}

\textbf{Methodological contributions} include the first multi-market machine learning analysis of UK BESS bidding and the development of a decision path clustering technique to uncover tactical sub-patterns.

\textbf{Limitations} include the focus on BMU-registered assets only, potential misclassification of battery fuel types, reliance on PN as a wholesale proxy, the exclusion of Capacity Market data, and the single-year study period that may not generalise to other years. Future research should incorporate revenue outcomes, geographic and asset characteristics, more granular wholesale data, and longer time horizons to assess strategic evolution.

\vspace{0.5cm}
\section*{Acknowledgements}
The author acknowledges the Energy Institute, University College London, and Professor Aiden O'sullivan for supporting this research.

\begin{thebibliography}{99}

\bibitem{eia2024a}
EIA (2024). \textit{U.S. Energy Information Administration -- Independent Statistics and Analysis}. Available at: \url{https://www.eia.gov/analysis/studies/electricity/batterystorage/pdf/battery_storage_2021}.

\bibitem{martins2021}
Martins, J. and Miles, J. (2021). A techno-economic assessment of battery business models in the UK electricity market. \textit{Energy Policy}, 148, p.111938. doi:10.1016/j.enpol.2020.111938.

\bibitem{bush2024}
Bush, J. (2024). \textit{GB BESS Outlook Q2 2024: Battery revenue stacking and dispatch optimization}. Modo Energy.

\bibitem{jennings2025a}
Jennings, Z. (2025a). \textit{GB BESS: Why top-performing batteries earn 30\% higher revenues}. Modo Energy.

\bibitem{jennings2025b}
Jennings, Z. (2025b). \textit{GB Battery energy storage revenues hit a two-year high in December 2024}. Modo Energy.

\bibitem{modo2024}
Modo Energy (2024). \textit{GB Methodology -- Indices}. Available at: \url{https://modoenergy.com/benchmarking/methodology/gb}.


\bibitem{elexon2025}
Elexon (2025). \textit{BSCP 15: BM Unit Registration -- Elexon Digital BSC}. Available at: \url{https://bscdocs.elexon.co.uk/bsc-procedures/bscp-15-bm-unit-registration}.

\bibitem{neso2025}
NESO (2025). \textit{Enduring Auction Capability (EAC)}. Available at: \url{https://www.neso.energy/industry-information/balancing-services/enduring-auction-capability-eac}.

\bibitem{chen2022}
Chen, Q. et al. (2022). The competition and equilibrium in power markets under decarbonization and decentralization. \textit{iEnergy}, 1(2), pp.188--203. doi:10.23919/ien.2022.0025.

\bibitem{merten2020}
Merten, M. et al. (2020). Bidding strategy for battery storage systems in the secondary control reserve market. \textit{Applied Energy}, 268, p.114951. doi:10.1016/j.apenergy.2020.114951.

\bibitem{rajah2024}
Rajah, I. et al. (2024). Degradation-Aware Derating of Lithium-Ion Battery Energy Storage Systems in the UK Power Market. \textit{Electronics}, 13(19), 3817. doi:10.3390/electronics13193817.

\bibitem{fan2025}
Fan, F., Nwobu, J. and Campos-Gaona, D. (2025). Co-Located Battery Energy Storage Optimisation for Dynamic Containment Under the UK Frequency Response Market Reforms. \textit{CSEE JPES}. doi:10.17775/cseejpes.2023.01210.

\bibitem{harri2022}
Aaltonen, H. et al. (2022). Bidding a Battery on Electricity Markets and Minimizing Battery Aging Costs: A Reinforcement Learning Approach. \textit{Energies}, 15(14), 4960. doi:10.3390/en15144960.

\bibitem{seward2022}
Seward, W., Qadrdan, M. and Jenkins, N. (2022). Revenue stacking for behind the meter battery storage. \textit{Electric Power Systems Research}, 211, 108292. doi:10.1016/j.epsr.2022.108292.

\bibitem{tang2022}
Tang, Q., Guo, H. and Chen, Q. (2022). Multi-Market Bidding Behavior Analysis of Energy Storage System Based on Inverse Reinforcement Learning. \textit{IEEE Trans. Power Syst.}, 37(6), pp.4819--4831. doi:10.1109/tpwrs.2022.3150518.

\bibitem{weaver2022}
Weaver, N. (2022). \textit{The evolution of the GB battery energy storage revenue stack}. Modo Energy.

\bibitem{wendel2023a}
Wendel Hortop (2023a). \textit{Dynamic Containment Low: why have prices increased so much?}. Modo Energy.

\bibitem{wendel2024}
Wendel Hortop (2024). \textit{Balancing Reserve: Are ESO concerns over battery performance justified?}. Modo Energy.

\end{thebibliography}

\end{document}